\section{Analytic foundations and characteristic formulas}
\label{app:analytic}

This appendix proves the analytic inputs stated in
Section~\ref{sec:analytic}. Throughout, $M$ is the standing closed,
connected, oriented four-manifold with $\operatorname{Ric}=3g$ and
strictly positive sectional curvature. Only that normalized geometry
and the conventions of Section~\ref{sec:preliminaries} are used; none
of the volume, Hessian, coupling, or classification conclusions enters.

\subsection{Elliptic facts, density, and Poisson solvability}
\label{app:analytic-preliminaries}

We use the following basic analytic facts on a closed smooth
four-manifold. The inclusion $H^1\subset L^4$ is continuous, and
$H^1\subset L^s$ is compact for $2\le s<4$. Smooth scalar and bundle
operators with Laplacian leading symbol satisfy the local
$W^{2,p}$ estimates for $1<p<\infty$, together with their iterated
Sobolev consequences. All operators below have these hypotheses,
and all nonlinear terms are put on the right side before applying
the estimates. For scalar operators we use the strong interior estimates,
Dirichlet solvability, and improvement of strong-solution integrability
in Gilbarg--Trudinger~\cite[Theorem~9.11, pp.~235--236;
Theorem~9.15 and Lemma~9.16, pp.~241--242]{GilbargTrudinger}.
The initial weak-to-strong passage for the nonlinear minimizer is
justified in the next subsection, rather than inferred from an estimate
for strong solutions.
For a weak bundle equation with $L^2$ data, write the operator in a smooth
local frame: its principal part is scalar, while its first-order coupling
applied to an $H^1$ section belongs to $L^2$. A cutoff reduces each
component to a scalar equation for $\Delta_g-C$ on a smooth ball,
with $C>0$. Strong zero-boundary solvability and coercive energy
uniqueness identify that component with its $W^{2,2}$ solution.
Once this initial regularity is established, componentwise interior
estimates and interpolation absorb the smooth first-order couplings;
differentiation gives the higher regularity used for Poisson solutions
and eigenfunctions.

Smooth functions are dense in $H^1(M)$: use local mollification in a
finite coordinate cover and a smooth partition of unity. The same
construction gives uniform and $H^1$ approximation for a Lipschitz
function, as used in the weak norm testing rule.

For completeness, every smooth mean-zero function $h$ has a unique smooth
mean-zero solution of $\Delta\phi=h$. On the mean-zero subspace of $H^1$,
the Dirichlet energy is coercive. Indeed, otherwise a sequence with
unit $L^2$ norm and gradients tending to zero has, by compactness in
$L^2$, a mean-zero limit with unit norm and zero weak gradient. On the
connected manifold that limit is constant, a contradiction. The
variational equation
\[
 \mathbf E\langle d\phi,d\psi\rangle=-\mathbf E(h\psi)
\]
therefore has a unique mean-zero weak solution by the Hilbert-space
representation theorem. The local regularity above makes it smooth.
The same energy identity proves uniqueness. This is the Poisson fact
used in Proposition~\ref{prop:weighted-variance}.

\subsection{The sharp Sobolev inequality}
\label{app:sharp-sobolev}

We give the normalized four-dimensional Einstein case of Ilias's
inequality~\cite[Theorem~3]{Ilias}. Its nonlinear Hessian proof belongs
to the Bochner approach to elliptic rigidity; compare
Bidaut-V\'eron--V\'eron~\cite{BidautVeronVeron}.

\begin{proof}[Proof of Proposition~\ref{prop:sharp-sobolev}]
Fix $2<s<4$ and minimize
\[
 Q_s(w)=
 \frac{\mathbf E(w^2+\tfrac12|\nabla w|^2)}
      {(\mathbf E|w|^s)^{2/s}}
\]
under $\mathbf E|w|^s=1$. A minimizing sequence is bounded in
$H^1$; weak compactness, compact inclusion into $L^s$, and lower
semicontinuity give a minimizer. Taking its absolute value does
not increase the energy, so choose $w\ge0$.
The ordinary Sobolev embedding makes the minimum $m_s$ positive,
and the constant test function gives $m_s\le1$.
The Euler equation is
\[
 -\Delta w+2w=2m_s w^{s-1}.
\]
After setting $u=m_s^{1/(s-2)}w$, it becomes
\[
 -\Delta u+2u=2u^{s-1}.
\]

We record the regularity needed to use a negative power of $u$.
We first pass from the weak minimizer to a strong solution. Choose a
smooth coordinate ball $B$ and a cutoff $\eta\in C_c^\infty(B)$.
For $z=\eta u$ the weak equation gives
\[
 (\Delta_g-2)z
 =-2\eta u^{s-1}+2\langle d\eta,du\rangle+(\Delta_g\eta)u
 =:F\in L^{r_0}(B),
 \qquad r_0=\min\{2,4/(s-1)\}>4/3.
\]
The strong Dirichlet existence theorem
\cite[Theorem~9.15]{GilbargTrudinger} gives a solution
$w_D\in W^{2,r_0}(B)\cap W^{1,r_0}_0(B)$ with this right side.
The domain and coefficients are smooth, the operator is uniformly
elliptic, and its zeroth-order coefficient is $-2$.
In dimension four, $r_0>4/3$ implies
$W^{2,r_0}\subset W^{1,q}$ for some $q>2$; compatibility of Sobolev
traces therefore gives $w_D\in H^1_0(B)$.
The difference $z-w_D\in H^1_0(B)$ solves
$(\Delta_g-2)(z-w_D)=0$ weakly. Testing by $z-w_D$ shows that it
vanishes. Thus $u$ is locally $W^{2,r_0}$, and the strong interior
estimates and integrability improvement now apply
\cite[Theorem~9.11 and Lemma~9.16]{GilbargTrudinger}.

Initially $u\in L^4$, while $u^{s-2}\in L^r$ with
$r=4/(s-2)>2$. If $u\in L^q$, the equation has right side in
$L^p$ with $1/p=1/q+1/r$. The $W^{2,p}$ estimate and Sobolev
embedding improve $1/q$ by a fixed positive amount smaller than
$1/2-1/r$. Initially $p=4/(s-1)>1$. Iteration reaches $p>2$,
giving boundedness; at a borderline exponent $p=2$, use any
finite larger $q$ before the next iteration.

After boundedness, the equation gives $u\in W^{2,p}$ for every finite
$p$. Since $s-1>1$, the derivative of $t^{s-1}$ is bounded on the
bounded nonnegative range of $u$. Differentiating the equation once,
and applying the same local regularity argument to each first
coordinate derivative, gives $u\in W^{3,p}$ for every finite $p$.
In particular, $u\in C^{2,\alpha}$ for some $\alpha>0$.
For $C\ge2$ we have
\[
 (-\Delta_g+C)u=(C-2)u+2u^{s-1}\ge0.
\]
The classical strong maximum principle gives $u>0$, since
$u\ge0$ and $u\not\equiv0$. A direct barrier proof uses a positive
exponential radial function on an annulus tangent to the zero set:
the weak maximum principle forces a nonzero inward derivative at the
tangency point, contradicting nonnegativity and differentiability.
Compactness now gives a positive minimum. The nonlinearity is smooth
on this positive range, so further differentiation and elliptic
regularity give smoothness of $u$.

Put $F=u^{-1}$. Its equation is
\[
 F\Delta F-2|\nabla F|^2+2F^2=2F^{4-s}.
\]
Since
$\operatorname{div}(\operatorname{Hess}_0F)
 =\frac34\,d(\Delta F+4F)$, differentiation gives
\[
 \begin{aligned}
 \operatorname{div}(F^{-3}\operatorname{Hess}_0F)
 &=\frac34F^{-4}
 \bigl[F\,d\Delta F+(\Delta F)dF
       -4\operatorname{Hess}F(\nabla F,\,\cdot\,)+4F\,dF\bigr]\\
 &=\frac32(4-s)F^{-1-s}\,dF .
 \end{aligned}
\]
Pairing with $dF$ and integrating yields
\[
 -\mathbf E F^{-3}|\operatorname{Hess}_0F|^2
 =\frac32(4-s)\mathbf E F^{-1-s}|dF|^2.
\]
Both sides have opposite signs, so $dF=0$. The positive constant
solution of the equation for $u$ is $u=1$. The normalization of
$w$ then gives $w=1$ and $m_s=1$.

Thus the claimed inequality holds with $4$ replaced by every
$s\in(2,4)$ and outer exponent $2/s$. For fixed $f\in H^1$,
the ordinary inclusion $H^1\subset L^4$ and dominated convergence
give $\|f\|_{L^s(d\nu)}\to\|f\|_{L^4(d\nu)}$ as $s\uparrow4$.
This proves the stated critical inequality.
\end{proof}

\subsection{Harmonic forms and the characteristic identities}
\label{app:characteristic-identities}

We derive the Euler and signature formulas using the complex of forms and
an explicit signed heat computation. This is the heat-supertrace method
of McKean--Singer~\cite{McKeanSinger1967}; general formulations and local
coefficient calculations appear in Gilkey~\cite[Theorems~1.7.6 and
4.8.16]{Gilkey1984}.

We first record the spectral heat construction for a smooth symmetric
Laplace-type operator $L=-\Delta+E$ on a Euclidean bundle over $M$. Adding a sufficiently
large constant gives a coercive quadratic form on $H^1$.
Its weak inverse is compact and self-adjoint on $L^2$ by
compact embedding, and elliptic regularity makes its
eigenfunctions smooth. The spectrum is discrete with finite
multiplicities and tends to infinity. Spectral heat evolution
is smoothing for every positive time: polynomial powers of
the shifted eigenvalue times $e^{-t\lambda}$ are bounded,
and elliptic estimates for powers of $L$ identify the
corresponding Sobolev bounds.

For completeness, these estimates also justify smooth kernel
and trace expansions. For a sufficiently large integer $k$,
evaluation of any prescribed finite derivative after
$(L+C)^{-k}$ is a bounded functional on $L^2$, uniformly in
its base point. Applied to an orthonormal eigenbasis, the
Riesz representation theorem gives a uniform bound for the
sum of the squared derivative values divided by
$(\lambda+C)^{2k}$. Cauchy--Schwarz and the exponential factor
then give uniform convergence of the spectral kernel series
and all its derivatives for $t$ in compact subintervals of
$(0,\infty)$. Integrating the diagonal gives the usual heat
trace. Uniqueness with prescribed initial point mass follows
by testing the kernel against the eigenfunctions.

\begin{lemma}[The form-complex Laplacians]
\label{lem:form-complex-laplacians}
Consider the elliptic complex
\[
 \Lambda^0\xrightarrow{\ d\ }\Lambda^1
       \xrightarrow{\ \sqrt2\,d^+\ }\Lambda^+,
 \qquad d^+=P_+d.
\]
Its three Laplacians are the Hodge Laplacians on the indicated
bundles. Written as $-\Delta+E$, their potentials are
\[
 E_0=0,\qquad E_1=3\operatorname{Id},\qquad
 E_+=4\operatorname{Id}-2T.
\]
Its alternating heat trace is $1+n_+$.
\end{lemma}

\begin{proof}
The middle complex Laplacian is
\[
 dd^*+2(d^+)^*d^+
 =dd^*+d^*d+d^*\!\star d.
\]
The last term is zero by $d^2=0$ and the formula for the
codifferential. On self-dual two-forms, Hodge star exchanges
$dd^*$ and $d^*d$, so $2P_+dd^*$ is their sum. This identifies
the last Laplacian; the scalar one is immediate. Their elliptic
principal symbols are consequently those of the metric Laplacian.

Commuting covariant derivatives in the one-form Hodge Laplacian
gives the Ricci endomorphism, equal to $3\operatorname{Id}$.
For the plus potential, write a self-dual form as
$\eta=\sum_s\eta_s I_s/\sqrt2$. Its Hodge energy is
$2\mathbf E|d^*\eta|^2$, which equals
\[
 \mathbf E\left|\sum_s I_s\nabla\eta_s\right|^2
 =\mathbf E\left(
      |\nabla\eta|^2+\frac12\langle C\nabla\eta,\nabla\eta\rangle
            \right),
\]
where $C=\sum I_s\otimes\Sigma_s$ now uses the spin-one
generators on the plus coefficients. The equality follows by
expanding the quaternion products. Integrating the $C$ term
and taking its skew derivative commutator gives the potential
\[
 -\frac12\sum_s(1+t_s)\Sigma_s^2.
\]
In the spin-one representation,
$\Sigma_s^2=4(e_se_s^{\mathsf T}-\operatorname{Id})$.
The potential is therefore $4\operatorname{Id}-2T$,
as asserted.

The scalar kernel consists of the constants. A harmonic one-form
satisfies
\[
 0=\mathbf E\bigl(|\nabla\eta|^2+3|\eta|^2\bigr),
\]
so the middle kernel vanishes. The last kernel has dimension
$n_+$. On each positive eigenspace of the complex Laplacian,
the complex differential commutes with the Laplacian, and its
adjoint divided by the eigenvalue is a contracting homotopy.
That finite-dimensional complex has zero alternating dimension.
Summing the heat factors over the positive spectrum leaves only
the kernel contribution $1+n_+$.
\end{proof}

\begin{lemma}[Local signed heat coefficient]
\label{lem:signed-heat-coefficient}
For the three operators in
Lemma~\ref{lem:form-complex-laplacians}, let $\operatorname{str}$
denote the trace on $\Lambda^0$ minus that on $\Lambda^1$ plus
that on $\Lambda^+$. Write $\Omega$ for their connection
curvatures. The constant local coefficient of their alternating
heat trace, including the factor $(4\pi)^{-2}$, is
\[
 \frac1{16\pi^2}
 \left(\frac12\operatorname{str}E^2
       +\frac1{12}\operatorname{str}\sum_{k,i}\Omega_{ki}^2\right).
\]
\end{lemma}

\begin{proof}
Fix the source point of the kernel. In normal coordinates
$z$ around it, use radially parallel orthonormal frames on
each of the three bundles and set
$\mu=\sqrt{\det g}$, $h=\mu^{-1/2}$. The ansatz
\[
 (4\pi t)^{-2}e^{-|z|^2/(4t)}
       \sum_{j\ge0}t^jU_j(z)
\]
gives $U_0=h\operatorname{Id}$ and the transport equations
\[
 \left(z\cdot\partial+j+\frac12z\cdot\partial\log\mu\right)U_j
 =-LU_{j-1}.
\]
Here the radial connection coefficient is zero.
The scalar part follows also from
$\Delta|z|^2=8+2z\cdot\partial\log\mu$.
Solving the radial equation gives
\[
 U_j(z)=-h(z)\int_0^1 s^{j-1}
                  \bigl(h^{-1}LU_{j-1}\bigr)(sz)\,ds.
\]
These transport solutions compute the actual local asymptotics.
Indeed, truncate at any large $N$ and apply a cutoff inside a
normal ball. The resulting kernel has initial delta mass,
and the transport equations leave a residual uniformly
$O(t^{N-2})$; derivatives of the cutoff give exponentially
small terms. The heat evolution of a metric connection
with bounded potential has a bounded sup norm on bounded
time intervals, by the maximum principle for the squared
norm. Duhamel integration therefore makes the difference
from the true kernel $O(t^{N-1})$. Compactness permits a
uniform normal-ball radius and uniform constants in the
source point. Taking $N$ large justifies the diagonal
coefficient calculation below.

In these frames write the connection as $\partial_i+A_i$.
Every $A_i$ is skew-symmetric and radial gauge gives
\[
 A_i(z)=\frac12z^k\Omega_{ki}(0)+O(|z|^2).
\]
For example, the exact radial-gauge identity
$A_i(z)=\int_0^1s z^k\Omega_{ki}(sz)\,ds$ follows by contracting
the curvature equation with $z^k$ and integrating along
the radius.
The signed ranks and potentials satisfy, at every point,
\[
 \operatorname{str}\operatorname{Id}=1-4+3=0,
 \qquad
 \operatorname{str}E=-12+\operatorname{tr}(4\operatorname{Id}-2T)=0.
\]
In coordinate form,
\[
 L=-g^{ij}\bigl[
   (\partial_i+A_i)(\partial_j+A_j)
       -\Gamma^k_{ij}(\partial_k+A_k)\bigr]+E.
\]
Apply this to $U_0=h\operatorname{Id}$. All purely scalar terms
have zero supertrace. Every term linear in $A_i$ or a
coordinate derivative of $A_i$ has zero ordinary trace in
each bundle, by skewness. The potential has zero supertrace
pointwise. Thus the exact surviving expression is
\[
 \operatorname{str}(LU_0)(z)
 =-h(z)g^{ij}(z)\operatorname{str}(A_i(z)A_j(z)).
\]
In particular, derivatives of the varying potential have
already disappeared before the radial integral is taken.
The transport formula now gives
\[
 \begin{aligned}
 \operatorname{str}U_1(z)
 &=h(z)\int_0^1 g^{ij}(sz)
              \operatorname{str}(A_i(sz)A_j(sz))\,ds\\
 &=\frac1{12}z^kz^\ell
       \operatorname{str}\sum_i\Omega_{ki}(0)\Omega_{\ell i}(0)
       +O(|z|^3).
 \end{aligned}
\]
The coefficient is
$\frac14\int_0^1s^2\,ds=\frac1{12}$.
Scalar and metric corrections do not change this quadratic term.

At the center, $A_i(0)=0$, $\Gamma(0)=0$, and
$\sum_i\partial_iA_i(0)=\frac12\sum_i\Omega_{ii}(0)=0$
as an endomorphism. Consequently
\[
 U_1(0)=c\operatorname{Id}-E(0),\qquad
 (LU_1)(0)=-\sum_i\partial_i^2U_1(0)+E(0)U_1(0),
\]
where $c$ is the same scalar for all three bundles.
Thus there are no additional connection terms at the center.
Taking the Laplacian of the signed quadratic expression gives
\[
 \operatorname{str}\sum_i\partial_i^2U_1(0)
 =\frac16\operatorname{str}\sum_{ki}\Omega_{ki}^2,
 \qquad
 \operatorname{str}(EU_1(0))=-\operatorname{str}E^2.
\]
Since the transport equation at the center gives
$U_2(0)=-\frac12(LU_1)(0)$, we obtain
\[
 \operatorname{str}U_2(0)
 =\frac12\operatorname{str}E^2
       +\frac1{12}\operatorname{str}\sum_{ki}\Omega_{ki}^2.
\]
The lower coefficients have zero signed diagonal trace.
Multiplication by the heat prefactor $(4\pi t)^{-2}$
now proves the assertion.
\end{proof}

\begin{proof}[Proof of Proposition~\ref{prop:characteristic-identities}]
By Hodge theory, harmonic forms represent de Rham cohomology. The
scalar kernel is one-dimensional, and the one-form kernel vanishes by
Lemma~\ref{lem:form-complex-laplacians}. Hodge star gives the corresponding
dimensions in degrees four and three. On harmonic two-forms, the
intersection form is positive definite on the self-dual part and negative
definite on the anti-self-dual part. Thus its signature is $n_+-n_-$,
and the alternating sum of the Betti numbers is $2+n_++n_-$. This
identifies the integers in Proposition~\ref{prop:characteristic-identities}
with the topological signature and Euler characteristic without using
simple connectivity.

The potential-square supertrace is
\[
 \operatorname{str}E^2
 =-\operatorname{tr}(9\operatorname{Id}_{\Lambda^1})
   +\operatorname{tr}\bigl((4\operatorname{Id}_{\Lambda^+}-2T)^2\bigr)
 =12+24v^2.
\]
On one-forms, the connection curvature has plus and minus
generators $I_s,J_s$. Orthogonality of the two-form bases,
together with $I_s^2=J_s^2=-\operatorname{Id}$, gives
\[
 \operatorname{tr}_{\Lambda^1}\sum_{ki}\Omega_{ki}^2
 =-4\sum_s\bigl((1+t_s)^2+(1+u_s)^2\bigr)
 =-24(1+v^2+b^2).
\]
On plus two-forms only the plus generators act, and each
speed-two spin-one generator has square trace $-8$.
Therefore
\[
 \operatorname{tr}_{\Lambda^+}\sum_{ki}\Omega_{ki}^2
 =-8\sum_s(1+t_s)^2=-24(1+2v^2).
\]
The scalar connection is flat. Hence
$\operatorname{str}\sum\Omega_{ki}^2=24(b^2-v^2)$, and
Lemma~\ref{lem:signed-heat-coefficient} gives the signed
constant density
\[
 \frac1{16\pi^2}(6+10v^2+2b^2).
\]
Integrate this against $d\operatorname{vol}_g$ and compare
with the constant alternating heat trace from
Lemma~\ref{lem:form-complex-laplacians}:
\[
 16\pi^2(1+n_+)
 =\operatorname{Vol}(M)\,
       \mathbf E(6+10v^2+2b^2).
\]
Reverse orientation to get the analogous formula with
$n_-,b,v$. Adding and subtracting the two formulas gives
exactly \eqref{eq:6}.

On anti-self-dual forms the potential is
$4\operatorname{Id}-2U$. If $U=0$, the energy of a harmonic
anti-self-dual form is the integral of
$|\nabla\eta|^2+4|\eta|^2$, so its kernel is zero.
The assertion for $T=0$ follows in the same way.
\end{proof}
